%% ma: L11-B27-0011
%% lop: 11
%% chuong: 7
%% bai: 27
%% dang: 11.27.8
%% loai: TLN
%% muc: VD
%% dap_an: "0,24"
%% nguon: "(NBV)-ĐỀ SỐ 6-MỖI NGÀY 1 ĐỀ THI 2025.pdf (D06, MỖI NGÀY 1 ĐỀ - Đề số 6), Phần III Câu 2 (THPT Triệu Sơn 3 - Thanh Hóa 2025)"
%% kien_thuc: "Bài 27 (thể tích khối chóp cụt đều), Bài 23 (đường thẳng vuông góc với mặt phẳng), Định lí Pythagore"
%% trang_thai: cho-duyet
%% ly_do: "Dạng toán mới đề xuất 11.27.8 (thể tích khối chóp cụt đều từ cạnh hai đáy và cạnh bên), chờ giáo viên duyệt"
%% ghi_chu: "Đáp án gốc 0,24 đúng (V≈236594 cm3). Bỏ hình: hình gốc chỉ là ảnh chụp chiếc sọt đan (ảnh đời sống, không chứa dữ kiện toán), mọi số liệu đã nằm trong lời đề"
\begin{ex}
Một sọt đựng đồ có dạng hình chóp cụt đều (hai đáy là hai hình vuông song song, các cạnh bên bằng nhau). Đáy và miệng sọt là các hình vuông có cạnh tương ứng bằng $80$ cm và $60$ cm. Cạnh bên của sọt dài $50$ cm. Tính thể tích của sọt theo đơn vị mét khối, lấy kết quả đến hàng phần trăm.
\shortans{0,24}
\loigiai{Đáy lớn $ABCD$ có cạnh $80$ cm, đáy nhỏ $A'B'C'D'$ có cạnh $60$ cm, cạnh bên $AA'=50$ cm. Nửa đường chéo của hai đáy chênh nhau $\dfrac{80\sqrt2-60\sqrt2}2=10\sqrt2$ cm, đó là khoảng cách $DH$ từ $D$ đến hình chiếu $H$ của $D'$ lên đáy lớn theo phương đường chéo.

Chiều cao $h=\sqrt{50^2-(10\sqrt2)^2}=10\sqrt{23}$ cm. Diện tích hai đáy $S_1=80^2=6400$, $S_2=60^2=3600$.

$V=\dfrac13h\big(S_1+\sqrt{S_1S_2}+S_2\big)=\dfrac13\cdot10\sqrt{23}\,(6400+4800+3600)\approx236\,594\ \text{cm}^3\approx0{,}24\ \text{m}^3$.}
\end{ex}
